\documentclass[12pt]{article}
\usepackage[a4paper,margin=2.0cm]{geometry} % Standard margins for journals
\usepackage{setspace} % Line spacing control
\usepackage{titling} % Title positioning
\usepackage{enumitem} % List customization
\usepackage{amsmath,amssymb,amsthm} % Mathematics packages
\usepackage{graphicx} % Graphics inclusion
\usepackage{float} % Float positioning control
\usepackage{xcolor} % Color support
\usepackage{physics} % Physics notation
\usepackage{wrapfig} % Text wrapping around figures
\usepackage{array} % Table alignment control
\usepackage{pgfplots}
\pgfplotsset{compat=1.18}
\usepackage{booktabs}
\usepackage{multirow}
\usepackage{colortbl}
\usepackage{caption} % Caption customization
\usepackage{tabularx} % Extended table functionality
\usepackage{booktabs} % Professional tables
\usepackage{titlesec} % Section formatting
\usepackage{hyperref} % Hyperlinks and references
\usepackage{subcaption}
\usepackage{animate} % Animation support
\usepackage[numbers]{natbib} % Bibliography management with numerical citations
\usepackage{listings} % Code listing
\usepackage{extarrows}
\usepackage{cancel} % For strikethrough in math mode
\usepackage{colortbl} % For colored table rows
% Header and footer settings
\usepackage{fancyhdr}
\pagestyle{fancy} 
\fancyhf{} % clear default header/footer
\renewcommand{\sectionmark}[1]{%
  \markboth{\thesection\ #1}{}%
}
\renewcommand{\subsectionmark}[1]{%
  \markright{\thesubsection\ #1}%
}

% Header content: odd pages show Section; even pages show Section -- Subsection (if available)
\fancyhead[C]{%
  \nouppercase{%
    \ifodd\value{page}%
      \leftmark
    \else
      \leftmark
      \ifx\rightmark\empty\else\ -- \rightmark\fi
    \fi
  }%
}
% 頁尾頁碼
\fancyfoot[C]{\thepage}

\setlength{\headheight}{28pt}
\addtolength{\topmargin}{-13pt}

% Code listing settings for C++
\lstset{
    language=C++,
    basicstyle=\ttfamily\scriptsize,
    keywordstyle=\color{blue}\bfseries,
    commentstyle=\color{gray}\itshape,
    stringstyle=\color{red},
    numbers=left,
    numberstyle=\tiny\color{gray},
    stepnumber=1,
    numbersep=8pt,
    showstringspaces=false,
    breaklines=true,
    frame=single,
    frameround=tttt,
    tabsize=2,
    captionpos=b,
    backgroundcolor=\color{white},
    xleftmargin=3em,
    framexleftmargin=2.5em,
    escapeinside={(*@}{@*)},
    morekeywords={nullptr}
}

% 1. List spacing settings
\setlist[enumerate]{itemsep=2pt, parsep=0pt, topsep=2pt}
\setlist[itemize]{itemsep=2pt, parsep=0pt, topsep=2pt}

% 2. Caption settings
\captionsetup{
    labelfont={small,bf},    % Label: small font + bold
    textfont={small},         % Text: small font
    skip=10pt                 % Space between caption and figure
}

% 3. Table spacing settings
\setlength{\tabcolsep}{8pt} % Column separation
\renewcommand{\arraystretch}{1.3} % Row height

% 4. Footnote settings
\usepackage[bottom,hang]{footmisc} % Footnotes at bottom with hanging indent
\setlength{\footnotesep}{10pt} % Space between footnotes
\setlength{\skip\footins}{12pt plus 5pt minus 2pt} % Space between text and footnotes
\setlength{\footnotemargin}{1.5em} % Space between footnote mark and text
\renewcommand{\footnoterule}{\vspace*{-3pt}\hrule width 0.4\columnwidth height 0.4pt\vspace*{3pt}}

% 5. Theorem style with upright text
\newtheoremstyle{upright}
  {6pt}{6pt}  % Space before and after
  {\normalfont}% Body font (upright, not italic)
  {0pt}{\bfseries}{.}{0.5em}{}

% 6. Define theorem environments
\theoremstyle{plain}  % 定理、引理用斜體
\newtheorem{theorem}{Theorem}[section]  
\newtheorem{lemma}{Lemma}[section]
\newtheorem{corollary}{Corollary}[section]

\theoremstyle{definition}  % 定義、例子、註解用正體
\newtheorem{definition}{Definition}[section]
\newtheorem{example}{Example}[section]
\newtheorem{remark}{Remark}[section]


% 7. Section numbering format (standard Arabic numerals)
\renewcommand{\thesection}{\arabic{section}}
\renewcommand{\thesubsection}{\thesection.\arabic{subsection}}
\renewcommand{\thesubsubsection}{\thesubsection.\arabic{subsubsection}}
\renewcommand{\theequation}{\thesection.\arabic{equation}}
\renewcommand{\thefigure}{\thesection.\arabic{figure}}
\renewcommand{\thetable}{\thesection.\arabic{table}}

% 8. Reset equation counter at each section
\counterwithin{equation}{section}
\counterwithin{figure}{section}
\counterwithin{table}{section}

% 9. Section spacing adjustments
\titlespacing*{\section}
{0pt}                    % Left margin
{12pt}                   % Space before
{6pt}                    % Space after

\titlespacing*{\subsection}
{0pt}                    % Left margin
{10pt}                   % Space before
{4pt}                    % Space after

\titlespacing*{\subsubsection}
{0pt}                    % Left margin
{8pt}                    % Space before
{3pt}                    % Space after

% 10. Hyperref setup (should be last)
\hypersetup{
    colorlinks=true,
    linkcolor=blue,
    citecolor=blue,
    urlcolor=blue,
    bookmarksnumbered=true
}

% 11. Animation toggle switch
% Set to \showanimationtrue when finalizing, \showanimationfalse for faster compilation
\newif\ifshowanimation
\showanimationfalse  % Change to \showanimationtrue when ready to include animations

% 12. Title settings
\setlength{\droptitle}{-0.5in} % Adjust title position
\title{CFD Assignment 3 : Lid-Driven Cavity Flow}
\author{Chen Peng Chung }
\date{\today}

\begin{document}
\setcounter{section}{0}
\maketitle 
\tableofcontents
\newpage  
\section {Introduction}
For Assignment 3, solving Lid-Driven Cavity Flow,
I have used four types of methods to solve this classic problem. Corresponding to the difficulties I encountered during the Re number increasing or numerical oscillation happening for the high order schemes, I can make a fast review below: 
\begin{enumerate}
  \item \ref{subsec:code4}: FDM Central Difference Scheme\\
  This method is Finite Difference Method with SIMPLE algorithm to solve the unsteady Navier-Stokes Equations. The disadvantage of this code is the accuracy in space is only 2nd order, and time accuracy is only 1st order, because the discrete equation in this code is scheme of forward difference in time and central difference in space (FTCS). 
  \item \ref{subsec:code6}: FVM Deferred Correction\\
   I would like to describe this problem from the figure below \ref{fig:deferred correction}. When I tried to compute at Re = 5000 with non-deferred correction FVM code \ref{subsec:code5}, I found that whatever how I reduce the relaxation factor, I still can't get the appropriate iteration process, even only one converged solution at all. 
   The solving method I took is that changed the structure of the discrete equation of momentum equations, i.e., moving the anti-diffusion term to the source term, and meaning is using old value to substitute those terms, leading to elements of coefficient matrix maintain the Convection Boundedness Criterion(CBC) condition.  
\ifshowanimation
\begin{figure}[H]
    \centering
    \begin{subfigure}[b]{0.49\textwidth}
        \centering
        \animategraphics[loop,controls,width=\textwidth]{10}{}{1}{37}
        \caption{}
    \end{subfigure}
    \hfill
    \begin{subfigure}[b]{0.49\textwidth}
        \centering
        \animategraphics[loop,controls,width=\textwidth]{10}{}{38}{75}
        \caption{}
    \end{subfigure}
    \caption{Lid-driven cavity flow at $Re = 5000$ on $81 \times 81$ grid. Comparison of (a) non-deferred and (b) deferred correction schemes. The deferred correction technique maintains the Convection Boundedness Criterion (CBC) and ensures convergence. (Click animation to play in Adobe Acrobat Reader)}
    \label{fig:cavity_animation_split}
\end{figure}
\else
\begin{figure}[H]
    \centering
    \begin{subfigure}[b]{0.49\textwidth}
        \centering
        \includegraphics[width=\textwidth]{37.png}
        \caption{Non-deferred correction failure (Re = 5000)}
    \end{subfigure}
    \hfill
    \begin{subfigure}[b]{0.49\textwidth}
        \centering
        \includegraphics[width=\textwidth]{75.png}
        \caption{Deferred correction success (Re = 5000)}
    \end{subfigure}
    \caption{Comparison of correction methods for Lid-Driven Cavity Flow}
    \label{fig:cavity_comparison}
\end{figure}
\fi
    \item \ref{subsec:code7}: FVM Deferred Correction with ADI+TDMA\\
The last problem is about how can I reduce the computational time, furthermore, choosing the most appropriate relaxation factor set including factors of momentum equations and pressure factor in SIMPLE algorithm. When I was just using Gauss-Seidel method to solve the pressure correction equation, I had to wait for the converged iteration final for 3 days, and when I took the ADI (Alternating-Direction Implicit technique) to rearrange the discrete form of the momentum equation and pressure Poisson equation, the time I used reduced to 10 minutes. 
    \item Relaxation factor settings\\
    After I realized that I have to increase the computational efficiency, I also took the numerical experiment more and more again to find the most appropriate relaxation factor set. Finally, results I observed are summarized below: 
    \begin{itemize} 
    \item When pressure relaxation factor is smaller than velocity relaxation factor, the solution must be diverged despite how small the relaxation factors are. 
    \item In some normal range (e.g., pressure relaxation factor = 0.3, velocity relaxation factor = 0.05), the difference of both is larger, the final error compared with the benchmark is smaller. 
    \item The most appropriate relaxation factor set I found is: relaxation factor for pressure is 0.9, and for velocity is 0.01. 
    \end{itemize}
\end{enumerate}
\section{Problem Descriptions}

% === Lid-Driven Cavity Flow Animation ===
\ifshowanimation
\begin{figure}[H]
    \centering
    \begin{subfigure}[b]{0.51\textwidth}
        \centering
        \includegraphics[width=\textwidth]{Problem.png}
        \caption{Problem Description (Re = 100)}
    \end{subfigure}
    \hfill
    \begin{subfigure}[b]{0.47\textwidth}
        \centering
        \animategraphics[loop,controls,width=0.9\textwidth]{10}{ezgif-frame-}{001}{051}
        \caption{Lid-Driven Cavity Animation  (Re = 5000)}
    \end{subfigure}
    \caption{Lid-Driven Cavity Flow Problem and Simulation (Click animation to play in Adobe Acrobat Reader)}
    \label{fig:cavity_animation}
\end{figure}
\else
\begin{figure}[H]
    \centering
    \begin{subfigure}[b]{0.51\textwidth}
        \centering
        \includegraphics[width=\textwidth]{Problem.png}
        \caption{Problem Description (Re = 100)}
    \end{subfigure}
    \hfill
    \begin{subfigure}[b]{0.47\textwidth}
        \centering
        \includegraphics[width=\textwidth]{ezgif-frame-025.jpg}
        \caption{Simulation Result (Re = 5000)}
    \end{subfigure}
    \caption{Lid-Driven Cavity Flow Problem and Simulation (Animation disabled for faster compilation)}
\end{figure}
\fi


\noindent Lid Driven Cavity Flow can be controlled by Navier-Stokes Equation, and continuity equation. Writing down the conservation form of Navier-Stokes Equations below: 
\begin{equation}
\begin{aligned}
\frac{\partial \rho U}{\partial x} + \frac{\partial \rho V}{\partial y} = 0 \\[1.5ex]
\frac{\partial \rho U}{\partial t} + \frac{\partial \rho UU}{\partial x} + \frac{\partial \rho VU}{\partial y} &= -\frac{\partial P}{\partial x} + \mu\left[\frac{\partial^2 U}{\partial x^2} + \frac{\partial^2 U}{\partial y^2}\right] \\[1.5ex]
\frac{\partial \rho V}{\partial t} + \frac{\partial \rho UV}{\partial x} + \frac{\partial \rho VV}{\partial y} &= -\frac{\partial P}{\partial y} + \mu\left[\frac{\partial^2 V}{\partial x^2} + \frac{\partial^2 V}{\partial y^2}\right]
\end{aligned}
\end{equation}
\begin{enumerate}[label=(\alph*)]
    \item Please compute the cavity flows with Reynolds number being 100, 1000 and 5000.  
    \item Please use mesh size at least 81x81 and 161x161.
    \item Please also compare your results with the benchmark solutions from Ghia et al. (1982) for u(y) at
x=0.5 and v(x) at y=0.5.
    \item Please compare the accuracy of the central difference, QUICK and MUSCL schemes.
\end{enumerate}
and I will show the analysis below to answer these problems.
\\[1.5ex]
{\color{blue}
\begin{enumerate}
    \item Streamline plots and velocity contours at Re = 100, 1000, 5000 with mesh size 81x81 and 161x161 for CDs, QUICK, MUSCL scheme.(36 figures)
    \item centerline velocity profiles $u(y)$ and $v(x)$ comparison with different Reynolds number. (12 figures)
    \item centerline velocity profiles comparison with different scheme. (6 figures)
    \item centerline velocity profiles comparison with mesh size (Grid 81x81, 113x113, 161x161, 227x227). (4 figures)
    \item Accuracy analysis and discussion.
  \end{enumerate}
}

\section{Introduction of Methodology Adopted}
\subsection{Discretize the Governing Equations}
\subsubsection{Staggered Grid: Problem and Solving Method}


\begin{definition}
\noindent Construct three types of mesh system : 
\begin{enumerate}
\label{def:staggered_grid}
\item primary point: The storage point of pressure or other scalar variables, and located at $(I,J)$.   
\item x-Velocity vice point: The storage point of x component of velocity, and located at $(i,J)$.   
\item y-Velocity vice point: The storage point of y component of velocity, and located at $(I,j)$.   
\end{enumerate}
\end{definition}

\begin{definition}
\noindent For solving flow field problem, the centers of Navier-Stokes equations were split into three parts :  
\begin{enumerate}
\label{def:center of the governing equation}
\item control volume center of Continuity Equation : primary point $(I,J)$
\item control volume center of x-momentum equation : x-Velocity vice point $(i,J)$
\item control volume center of y-momentum equation : y-Velocity vice point $(I,j)$
\end{enumerate}
\end{definition}

\begin{figure}[H]
    \centering
    \begin{minipage}[b]{0.48\textwidth}
        \centering
        \includegraphics[width=\textwidth]{76.png}
        \\[0.5ex]
        \caption{\footnotesize definition of primary point and vice point} 
    \end{minipage}\hfill
    \begin{minipage}[b]{0.48\textwidth}
        \centering
        \includegraphics[width=\textwidth]{77.png}
        \\[0.5ex]
        \caption{\footnotesize control volume center setting}
    \end{minipage}
    \caption{staggered grid system}
    \label{fig:staggered-grid}
\end{figure}

\noindent Why do we must use staggered grids  to deal with flow field problem? 
The problem generated from difference scheme of the momentum equation, especially the discrete form of pressure gradient term.\\
\begin{equation}
  \label{eq:problem}
  -\frac{1}{\rho}\oint_{\partial \Omega} dy \cdot p \approx 
  -\frac{1}{\rho} (p_{e} - p_{w})    \approx
  -\frac{1}{2\rho} (p_{E} - p_{W}) 
\end{equation}
so from equation \eqref{eq:problem}, if we take the pressure gradient as central difference scheme, we will face the difficulty that difference distance of it is "2-$\delta x$". However, in SIMPLE algorithm, i.e., 
\[
\text{SIMPLE algorithm is a type of pressure and velocity coupling correction method.}
\]
we give an initial pressure predicted value at first, and solve the predicted form of momentum equation to get velocity predicted value. Furthermore, we use pressure correction equation to update pressure corrected value, and then update the next step pressure predicted value.
\begin{figure}[H]
    \centering
    \begin{minipage}[b]{0.38\textwidth}
        \centering
        \includegraphics[width=\textwidth]{78.png}
        \\[0.5ex]
        {\footnotesize (a) coupling correction relationship} 
    \end{minipage}\hspace{5em}
    \begin{minipage}[b]{0.35\textwidth}
        \centering
        \includegraphics[width=\textwidth]{79.png}
        \\[0.5ex]
        {\footnotesize (b) checkerboard pressure field}
    \end{minipage}
    \caption{discussion of reason why we adopted staggered grids setting}
    \label{fig:discussion}
\end{figure}
As figure \ref{fig:discussion} shows, we can find that the velocity and pressure fields are corrected in a coupled manner in the SIMPLE algorithm. If the pressure predicted value adds a sawtooth wave in the process of updating, the velocity generated from $2\text{-}\delta x$ pressure gradient will not correct the pressure correction value, and correct pressure predicted value at the next step. The result of the iteration is that the non-physical error will be maintained to the iteration final. 
\begin{center}
Velocity generated from $1\text{-}\delta x$ pressure gradient\\
can modify the non-physical phenomenon of the pressure predicted value updating.
\end{center}
So this is the reason why we must use staggered grid setting, because we have to create the velocity predicted value generated from $1\text{-}\delta x$ pressure gradient. If the velocity field is stored at its vice point, the corresponding pressure gradient can directly adopt the form: $\frac{A_{e}p_{e} - A_{w}p_{w}}{\rho}$, and then avoid the checkerboard pressure field generated at this moment. 
\subsubsection{Discrete Form of the Continuity Equation}
\noindent The integration form of the Continuity Equation is : 
\begin{equation}\begin{aligned}
   \oint_{\partial \Omega} d\vec{a} \cdot \rho \vec{u} = 0 \\[1.5ex]
\end{aligned}\end{equation}
\begin{itemize}
  \item center of the control volume : primary point $(I,J)$, as the definition \ref{def:center of the governing equation} states.
  \item For incompressible flow and steady state flow, the density $\rho$ is constant in space and time, so we can simplify the equation to :$\oint_{\partial \Omega}d\vec{a} \cdot \vec{u} = 0$ 
\end{itemize}
\noindent The discrete form of the Continuity Integration Equation is :
\begin{equation}
  \label{eq:discrete_continuity_equation}
  (u_{e}A_{e} - u_{w}A_{w}) + (v_{n}A_{n} - v_{s}A_{s}) = 0
\end{equation} 
Furthermore, because we adopted the staggered grids setting, the variables $v_{n}, v_{s}$ stored at the north and south side of the control volume, and variables $u_{e} , u_{w}$ also stored at the east and west side.
\begin{figure}[H]
  \centering
  \includegraphics[width=0.7\textwidth]{81.png}
  \caption{linear combination coefficients of the discrete Continuity Equation}
  \label{fig:discrete_continuity_equation}
\end{figure}
\textbf{General discrete form of the Continuity Equation:}
\begin{equation}
  F_{e}^{P} - F_{w}^{P} + F_{n}^{P} - F_{s}^{P} = 0
  \label{eq:general_C_E}
\end{equation} 
where $F_{e}^{P} = u_{e}A_{e}$, $F_{w}^{P} = u_{w}A_{w}$, $F_{n}^{P} = v_{n}A_{n}$, $F_{s}^{P} = v_{s}A_{s}$ as the figure \ref{fig:discrete_continuity_equation} shown.
\subsubsection{General Discrete Form of the Momentum Integration Equations} 
For the title "General", meaning is that the center value of the each side of the control volume is substituted by the face variables first, and high order scheme \footnote{high order schemes in this assignment 
are Central Difference scheme, QUICK, and TVD-MUSCL scheme, talking about them at the next subsubsection.} will discuss in the next section.
The Derivation process below discusses the x-momentum integration equation as an example.
For incompressible flow and steady state, \textbf{x-momentum integration equation} is : 
\begin{equation}
\label{eq:x-momentum_integration_equation}
\underbrace{\oint_{\partial\Omega} dy\, u \cdot u + dx\, v \cdot u}_{\text{Convection term}} = \underbrace{\int_{\Omega} d^3r \left( -\frac{1}{\rho} \frac{\partial p}{\partial x} \right)}_{\text{Pressure gradient term}} + \underbrace{\frac{1}{R_e} \oint_{\partial\Omega} \left(dy\, \frac{\partial u}{\partial x} + dx\, \frac{\partial u}{\partial y}\right)}_{\text{Diffusion term}}
\end{equation}
where $\Omega$ is the control volume, and $\partial \Omega$ is the close surface of the control volume.
\begin{itemize}
  \item center of the control volume : x-Velocity face point $(i,J)$
\end{itemize}
\begin{figure}[H]
  \centering
  \includegraphics[width=0.7\textwidth]{82.png}
  \caption{coefficients of the general discrete form of the x-momentum integration equation}
  \label{fig:discrete_x_momentum_equation}
\end{figure}
For discretizing the equation \eqref{eq:x-momentum_integration_equation}, we start from discretizing each part of the equation :
\begin{enumerate}
  \item pressure gradient term : the method is that I took the 2nd order Central Difference of the first order partial derivative of pressure below.
  \begin{equation}
  \label{eq:pressure_CDS}
  -\frac{1}{\rho}\oint_{\partial \Omega} dy \cdot p \approx 
  -\frac{1}{\rho} (p_{e} - p_{w}) = p_{i,J} - p_{i-1,J}  
\end{equation}
  \item momentum diffusion term : the key point about discretizing this part is using 2nd order Central Difference to approximate the first order partial derivative of the u-velocity located at the center of each side of the control volume.   
  \begin{equation}
    \frac{1}{Re} \int_{e} dy \cdot \frac{\partial u}{\partial x} 
    \xlongequal{\text{Central Difference}} \frac{A_{e}}{Re\delta x}\left(u_{E} - u_{P}\right)
  \end{equation}
  where $\frac{A_{e}}{Re \delta x}$ is often defined as the symbol $D_{e}$, called it as "discrete diffusion coefficient". 
  \item momentum convection term : 
  the difficulties of discretizing the convection term is how to decide the x-velocity values at the east and west of the control volume, i.e., $u_{e}^{old} , u_{w}^{old}$, and decide the y-velocity at the north and south face of the control volume, i.e., $v_{n}^{old} , v_{s}^{old}$  
  ,as the figure \ref{fig:discrete_x_momentum_equation} shown. 
  {\color{red}
  \begin{center}
    Using 2nd order Central Difference Scheme to approximate the variables at each face of the control volume of the x-momentum equation.\footnote{The approximating method of the $u_{e},u_{w},v_{n},v_{s}$ is independent of the equation that the variables maintain, so I call the method as "Central Difference" for all variables.} 
  \end{center}
  }
  \begin{enumerate}[label=(\arabic*)]
    \item x-velocity at the east face of the control volume:
    \begin{equation}
      u_{e} \overset{\text{CDs}}{=} \frac{u_{E} + u_{P}}{2} = \frac{u_{i+1,J} + u_{i,J}}{2}
    \end{equation}
    \item x-velocity at the west face of the control volume:
    \begin{equation}
      u_{w} \overset{\text{CDs}}{=} \frac{u_{P} + u_{W}}{2} = \frac{u_{i,J} + u_{i-1,J}}{2}
    \end{equation}
    \item y-velocity at the north face of the control volume:
    \begin{equation}
      v_{n} \overset{\text{CDs}}{=} \frac{v_{i,J+1} + v_{i-1,J+1}}{2}
    \end{equation}
    \item y-velocity at the south face of the control volume:
    \begin{equation}
      v_{s} \overset{\text{CDs}}{=} \frac{v_{i,J} + v_{i-1,J}}{2}
    \end{equation}
  \end{enumerate}
  \begin{equation}
    \label{eq:convection_term_CDs}
    \oint_{\partial \Omega} dy \, u \cdot u + dx \, v \cdot u \approx 
    (u_{e}A_{e}u_{e}^{old} - u_{w}A_{w}u_{w}^{old}) + (v_{n}A_{n}u_{n}^{old} - v_{s}A_{s}u_{s}^{old})
  \end{equation}
  where $u_{e}^{old}, u_{w}^{old}, v_{n}^{old}, v_{s}^{old}$ are generated from last iteration step. This equation is the general discrete form of the momentum convection term.
\end{enumerate}
\noindent Combining equations \eqref{eq:pressure_CDS}, \eqref{eq:convection_term_CDs}, and diffusion term together, we can get \textbf{the general discrete form of the x-momentum integration equation below} :
\begin{align}
&F_e^u u_e - F_w^u u_w + F_n^u u_n - F_s^u u_s = \frac{1}{\rho} \frac{-(p_{I,J} - p_{I-1,J})}{\delta x} \delta x \delta y \notag \\
&\quad -u_P (D_w + D_e + D_s + D_n) + u_W (D_w) + u_E (D_e) + u_S (D_s) + u_N (D_n)
\end{align}
\subsubsection{High Order Scheme of the Momentum Integration Equations} 
\noindent "High order scheme" means that take high order approximation method to center value of each side of the control volume, so we can write down the general form of the approximation about them with flux limiter function :
\begin{equation}\begin{aligned}
  u_{f} &\xlongequal{\text{general approximation}}{=} D(r_{e})u_{D} + U(r_{f})u_{U} + UU(r_{f})u_{UU} \\[1ex]
        &\xlongequal{\text{CDs approximation}}{=} 0.5u_{D} + 0.5u_{U} + 0u_{UU} \\[1ex]
        &\xlongequal{\text{QUICK approximation}}{=} \frac{3}{8}u_{D} + \frac{6}{8}u_{U} + \frac{-1}{8}u_{UU} \\[1ex]
        &\xlongequal{\text{TVD-MUSCL approximation}}{=} d(r_{e})u_{D} + c(r_{f})u_{U} + uu(r_{f})u_{UU}
\end{aligned}\end{equation}
\vspace{-0.5em}
\begin{figure}[H]
  \centering
  \includegraphics[width=0.3\textwidth]{83.png}
  \caption{\footnotesize combination coefficients of TVD-MUSCL scheme approximation}
  \label{fig:TVD_MUSCL}
\end{figure}
\vspace{-0.8em}
For incompressible flow and steady state flow, \textbf{high order scheme of the x-momentum integration equation} : 
\begin{equation}\label{eq:general x-mom}\begin{aligned}
    a_{P}u_{P} 
    &+a_{EE}u_{EE} + a_{E}u_{E} + a_{W}u_{W} + a_{WW}u_{WW}\\
    &+a_{NN}u_{NN} + a_{N}u_{N} + a_{S}u_{S} + a_{SS}u_{SS}\\
    &=\frac{1}{\rho}\frac{-\left(p_{P} - p_{w}\right)}{\delta x} \delta x \delta y
\end{aligned}\end{equation}
\noindent Below we will set the different linear combination coefficients of the equation \eqref{eq:general x-mom}
for different high order scheme,i.e., QUICK and TVD-MUSCL scheme.\\
\noindent Discussion : 
(1) For QUICK scheme, the linear combination coefficients are :
\begin{equation}
\left\{
\begin{aligned}
a_{P} &=-(a_{EE}+a_{E}+a_{W}+a_{WW}+a_{NN}+a_{N}+a_{S}+a_{SS})+F_{e}^{u}-F_{w}^{u}+F_{n}^{u}-F_{s}^{u} \\[1ex]
a_{WW} &= 0 + (\alpha_{w}\frac{-1}{8})(-F_{w}^{u}) \\[1ex]
a_{W} &= (\alpha_{e}\frac{-1}{8})F_{e}^{u} + ((1-\alpha_{w})\frac{3}{8} +\alpha_{w}\frac{6}{8})(-F_{w}^{u}) -D_{e}\\[1ex]
a_{E} &= ((1-\alpha_{e})\frac{6}{8} +\alpha_{e}\frac{3}{8}) F_{e}^{u} + ((1-\alpha_{w})\frac{-1}{8})(-F_{w}^{u}) -D_{w}\\[1ex]
a_{EE} &= ((1-\alpha_{e})\frac{-1}{8})F_{e}^{u} + 0\\[1ex]
a_{SS} &= 0 + (\alpha_{s}\frac{-1}{8})(-F_{s}^{u})\\[1ex] 
a_{S} &= (\alpha_{n}\frac{-1}{8})F_{n}^{u} + ((1-\alpha_{s})\frac{3}{8} +\alpha_{s}\frac{6}{8})(-F_{s}^{u}) -D_{n}\\[1ex]
a_{N} &= ((1-\alpha_{n})\frac{6}{8} +\alpha_{n}\frac{3}{8})F_{n}^{u} +  ((1-\alpha_{s})\frac{-1}{8})(-F_{s}^{u}) -D_{s}\\[1ex]
a_{NN} &= ((1-\alpha_{n})\frac{-1}{8})F_{n}^{u} + 0
\label{eq:1}
\end{aligned}
\right.
\end{equation}
\vspace{-0.5em}
\noindent As the figure shown : 
\begin{figure}[H]
    \centering
    \includegraphics[width=\textwidth]{84.png}
    \caption{\footnotesize coefficients setting of QUICK scheme}
    \label{fig:QUICK_scheme}
\end{figure}

(2) For TVD-MUSCL scheme, the linear combination coefficients are :
\begin{equation}\left\{\begin{aligned}
a_{P} &=-(a_{EE}+a_{E}+a_{W}+a_{WW}+a_{NN}+a_{N}+a_{S}+a_{SS})\\[1.5ex]
&\quad +\left(D(r_{w})+U(r_{w})+UU(r_{w})\right)(-F_{w}^{u})\\[1.5ex]
&\quad +\left(D(r_{e})+U(r_{e})+UU(r_{e})\right)F_{e}^{u}\\[1.5ex]
&\quad +\left(D(r_{s})+U(r_{s})+UU(r_{s})\right)(-F_{s}^{u})\\[1.5ex]
&\quad +\left(D(r_{n})+U(r_{n})+UU(r_{n})\right)F_{n}^{u}\\[1.5ex]
a_{WW} &= 0 + (\alpha_{w}UU(r_{w}))(-F_{w}^{u}) \\[1.5ex]
a_{W} &= (\alpha_{e}UU(r_{e}))F_{e}^{u} + ((1-\alpha_{w})D(r_{w})+\alpha_{w}U(r_{w}))(-F_{w}^{u}) - D_{w}\\[1.5ex]
a_{E} &= ((1-\alpha_{e})U(r_{e})+\alpha_{e}D(r_{e}))F_{e}^{u} + ((1-\alpha_{w})UU(r_{w}))(-F_{w}^{u}) - D_{e}\\[1.5ex]
a_{EE} &= ((1-\alpha_{e})UU(r_{e}))F_{e}^{u} + 0 \\[1.5ex]
a_{SS} &= 0 + (\alpha_{s}UU(r_{s}))(-F_{s}^{u})\\[1.5ex]
a_{S} &= (\alpha_{n}UU(r_{n}))F_{n}^{u} + ((1-\alpha_{s})D(r_{s})+\alpha_{s}U(r_{s}))(-F_{s}^{u}) - D_{s}\\[1.5ex]
a_{N} &= ((1-\alpha_{n})U(r_{n})+\alpha_{n}D(r_{n}))F_{n}^{u} + ((1-\alpha_{s})UU(r_{s}))(-F_{s}^{u}) - D_{s}\\[1.5ex]
a_{NN} &= ((1-\alpha_{n})UU(r_{n}))F_{n}^{u} + 0 \\[1.5ex]
\label{eq:2}
\end{aligned}\right.\end{equation}
\vspace{-0.5em}
\noindent As the figure shown : 
\begin{figure}[H]
    \centering
    \includegraphics[width=\textwidth]{85.png}
    \caption{\footnotesize coefficients setting of TVD-MUSCL scheme}
    \label{fig:TVD_MUSCL_scheme}
\end{figure}


\vspace{0.5em}
\noindent \textbf{Remark:} The corresponding C++ implementation of the TVD-MUSCL scheme coefficients in equation \eqref{eq:2} can be found in the non-deferred correction code (see Appendix \ref{subsec:code5}). Below we present the key code snippets that demonstrate how these theoretical coefficients are computed in practice.

\begin{remark}[Implementation of TVD-MUSCL Coefficients]
\label{rem:TVD_implementation}
The implementation follows a systematic approach for each control volume face (west, east, south, north). For each face, we compute the contributions to neighboring coefficients ($a_{WW}, a_{W}, a_{E}, a_{EE}, a_{SS}, a_{S}, a_{N}, a_{NN}$) and the source term based on the flux limiter functions $D(r)$, $U(r)$, and $UU(r)$.
\end{remark}

\paragraph{West Face Implementation (lines 405--409):}
The west face contributes to coefficients $a_{WW}$, $a_{W}$, $a_{E}$, and $a_{P}$ according to equation \eqref{eq:2}:
\begin{lstlisting}[language=C++, caption={West face coefficient computation}, label={lst:west_coeff}]
a_WW[i-1][J-1] += (alpha_w_u[i-1][J-1]*UU(r_w_u[i-1][J-1]))*(-F_w_u[i-1][J-1]) ;
a_W[i-1][J-1] += ((1-alpha_w_u[i-1][J-1])*D(r_w_u[i-1][J-1])+alpha_w_u[i-1][J-1]*U(r_w_u[i-1][J-1]))*(-F_w_u[i-1][J-1]) - D_w ;
a_E[i-1][J-1] += ((1-alpha_w_u[i-1][J-1])*UU(r_w_u[i-1][J-1]))*(-F_w_u[i-1][J-1]);
a_P[i-1][J-1] += ((1-alpha_w_u[i-1][J-1])*U(r_w_u[i-1][J-1])+alpha_w_u[i-1][J-1]*D(r_w_u[i-1][J-1]))*(-F_w_u[i-1][J-1])+D_w ;
S1[i-1][J-1] += -(-p_star[i-2][J-1])*dy/rho ;
\end{lstlisting}

\paragraph{East Face Implementation (lines 417--421):}
The east face contributes to coefficients $a_{EE}$, $a_{E}$, $a_{W}$, and $a_{P}$:
\begin{lstlisting}[language=C++, caption={East face coefficient computation}, label={lst:east_coeff}]
a_EE[i-1][J-1] += (1-alpha_e_u[i-1][J-1])*UU(r_e_u[i-1][J-1])*(F_e_u[i-1][J-1]) ;
a_E[i-1][J-1] += (alpha_e_u[i-1][J-1]*D(r_e_u[i-1][J-1])+(1-alpha_e_u[i-1][J-1])*U(r_e_u[i-1][J-1]))*(F_e_u[i-1][J-1]) - D_e ;
a_W[i-1][J-1] += (alpha_e_u[i-1][J-1])*UU(r_e_u[i-1][J-1])*(F_e_u[i-1][J-1]) ;
a_P[i-1][J-1] += ((1-alpha_e_u[i-1][J-1])*D(r_e_u[i-1][J-1])+alpha_e_u[i-1][J-1]*U(r_e_u[i-1][J-1]))*(F_e_u[i-1][J-1])+D_e;   
S1[i-1][J-1] += -(p_star[i-1][J-1])*dy/rho ;
\end{lstlisting}

\paragraph{South Face Implementation (lines 445--448):}
The south face contributes to coefficients $a_{SS}$, $a_{S}$, $a_{N}$, and $a_{P}$:
\begin{lstlisting}[language=C++, caption={South face coefficient computation}, label={lst:south_coeff}]
a_SS[i-1][J-1] += (alpha_s_u[i-1][J-1]*UU(r_s_u[i-1][J-1]))*(-F_s_u[i-1][J-1]) ;
a_S[i-1][J-1] += ((1-alpha_s_u[i-1][J-1])*D(r_s_u[i-1][J-1])+(alpha_s_u[i-1][J-1])*U(r_s_u[i-1][J-1]))*(-F_s_u[i-1][J-1])-D_s ;
a_N[i-1][J-1] += ((1-alpha_s_u[i-1][J-1])*UU(r_s_u[i-1][J-1]))*(-F_s_u[i-1][J-1]) ;
a_P[i-1][J-1] += ((1-alpha_s_u[i-1][J-1])*U(r_s_u[i-1][J-1])+alpha_s_u[i-1][J-1]*D(r_s_u[i-1][J-1]))*(-F_s_u[i-1][J-1])+D_s ; 
\end{lstlisting}

\paragraph{North Face Implementation (lines 473--476):}
The north face contributes to coefficients $a_{NN}$, $a_{N}$, $a_{S}$, and $a_{P}$:
\begin{lstlisting}[language=C++, caption={North face coefficient computation}, label={lst:north_coeff}]
a_NN[i-1][J-1] += (1-alpha_n_u[i-1][J-1])*UU(r_n_u[i-1][J-1])*(F_n_u[i-1][J-1]) ;
a_N[i-1][J-1] += (alpha_n_u[i-1][J-1]*D(r_n_u[i-1][J-1])+(1-alpha_n_u[i-1][J-1])*U(r_n_u[i-1][J-1]))*(F_n_u[i-1][J-1]) - D_n ;
a_S[i-1][J-1] += (alpha_n_u[i-1][J-1])*UU(r_n_u[i-1][J-1])*(F_n_u[i-1][J-1]) ;
a_P[i-1][J-1] += ((1-alpha_n_u[i-1][J-1])*D(r_n_u[i-1][J-1])+alpha_n_u[i-1][J-1]*U(r_n_u[i-1][J-1]))*(F_n_u[i-1][J-1])+D_n ;
\end{lstlisting}

\vspace{0.5em}
Let's go back to the equation \eqref{eq:general x-mom}, 
we only need to change the coefficients setting (i.e., equations \eqref{eq:1} and \eqref{eq:2})
to get the different high order discrete scheme of the x-momentum integration equation. 
Similarly, we can get the high order scheme of y-momentum integration equation with the same method, and the thing you have to watch out is that coefficients setting is very complicated, in the next section, I will simplify the form of the discrete equation using deferred correction. 
\subsection{SIMPLE Algorithm}
\subsubsection{Predicted-Correction Method}
\noindent In predicted-correction method, split the variables, i.e., $u$-velocity, $v$-velocity, pressure field, into three parts, i.e., correct value, correction value and predicted value. 
\begin{equation}\begin{aligned}
    u_{i,J} &= u_{i,J}^{\prime} + u_{i,J}^{\star}\\[1.5ex] 
    v_{I,j} &= v_{I,j}^{\prime} + v_{I,j}^{\star}\\[1.5ex] 
    p_{I,J} &= p_{I,J}^{\prime} + p_{I,J}^{\star}
\end{aligned}\end{equation}
So in this method, we have increased the  number of variables from 3 to 9. 

\begin{center}\begin{minipage}[c]{0.2\textwidth}
    \textbf{3 variables:}
    \[
    \left\{
    \begin{aligned}
        &u_{i,J} \\[1.5ex]
        &v_{I,j} \\[1.5ex]
        &p_{I,J}
    \end{aligned}
    \right.
    \]
\end{minipage}
$\xrightarrow{\text{Predicted-Correction Method}}$
\begin{minipage}[c]{0.2\textwidth}
    \textbf{9 variables:}
    \[
    \begin{bmatrix}
    u_{i,J} & u_{i,J}^{*} & u_{i,J}^{\prime} \\
    v_{I,j} & v_{I,j}^{*} & v_{I,j}^{\prime} \\
    p_{I,J} & p_{I,J}^{*} & p_{I,J}^{\prime} \\
    \end{bmatrix}
    \]
\end{minipage}\end{center}
{\color{blue}Problem Identification: }
\noindent In Navier-Stokes equations, pressure and velocity fields are coupled with each other; therefore, no variable can be solved by itself alone in the momentum equation system. 
The solving method for the problem above is that when I give a pressure field, I can get the corresponding velocity from the momentum equations.
The implication is that the correct pressure can yield the correct velocity. However, how can I define a correct pressure? We have another governing equation, the Continuity Equation,
$$A_{e}u_{i+1,J} - A_{w}u_{i,J} + A_{n}v_{I,j+1} - A_{s}v_{I,j} = 0$$, so we hope that this equation can be equivalent to another form for pressure maintenance. 
{\color{red}Solving Method: }
We set the variables into two parts, i.e., predicted value and correction value. 
First, given the pressure predicted value $p_{I,J}^{\star}$, we can obtain the corresponding velocity predicted value. 
Second, use the velocity I get to go back solve the presswure corection and correct the pressure in the next step.
so in this algorithm, we can find a thing is pressure and velocity are corrected in an iterative way using predicted value and correction value.
\subsubsection{Patankar's Omission in SIMPLE Algorithm}
\subsubsection{Derivatation of Pressure Correction Equation}
\newpage
\section{Results and Discussion}
\subsection{Vortex Structure Preview}
\noindent  Three figures below shows that the vortex structure of Lid-Driven Cavity Flow at different Re numbers(Re = 100, 1000, 5000)
\begin{figure}[H]
\centering
\begin{subfigure}[b]{0.48\textwidth}
  \centering
  \includegraphics[width=\textwidth]{P_1.png}
  \caption{}
  \label{fig:Preview_1}
\end{subfigure}\hfill
\begin{subfigure}[b]{0.48\textwidth}
  \centering
  \includegraphics[width=\textwidth]{P_2.png}
  \caption{}
  \label{fig:Preview_2}
\end{subfigure}
\caption{Vortex structures at Re = 100, 1000, (a) Re = 100, (b)Re = 1000}
\label{fig:vortex_preview}
\end{figure}
\noindent As the figure \ref{fig:vortex_preview} shown, we can see the fact that when Re number increases form 100 to 1000, despite of moving to the center of the domain for primary vortex, two second vortex structure at bottom-left corner and botto-right corner is gradually growing.
\begin{figure}[H]
  \centering
  \includegraphics[width=0.8\textwidth]{P_3.png}
  \caption{Vortex structure at Re = 5000}
  \label{fig:Preview_3}
\end{figure}
\noindent As shown in Figure \ref{fig:Preview_3}, when the Reynolds number equals 5000, a third secondary vortex appears at the top-left of the flow field. This phenomenon indicates that the vortex structure becomes more complicated as the inertia of the flow field increases.
(1)Red domain is the primary votex, (2)blue domain is the second vortex at top-left corner, (3)green domain is the second vortex at bottom-left corner of the flow field
,and(4)light-blue domain is the second vortex at bottom-right corner of the flow field.
\newpage
\subsection{Streamline Plots}
\noindent This section presents streamline patterns for three numerical schemes (Central Difference Scheme, QUICK, and TVD-MUSCL) across different Reynolds numbers (Re = 100, 1000, 5000) and grid resolutions (81×81 and 161×161).
\subsubsection{Arrangement of Streamline Patterns}
\noindent The below figure shows the arrangement method of my streamline plots results, split them into 81×81 and 161×161 parts at first, and separate figures into three parts based on 3 types of Re numbers. 
\begin{figure}[H]
  \centering
  \includegraphics[width=0.8\textwidth]{Arrangement1.jpeg}
  \caption{Arrangement method of streamline plots: organized by grid resolution (81×81 and 161×161) and Reynolds number (Re = 100, 1000, 5000).}
  \label{fig:arrangement}
\end{figure}

\subsubsection{Central Difference Scheme (CDs)}
\begin{figure}[H]
    \centering
    \begin{subfigure}[b]{0.42\textwidth}
        \includegraphics[width=\textwidth]{S_1_1_1.png}
        \caption{}
    \end{subfigure}
    \hfill
    \begin{subfigure}[b]{0.42\textwidth}
        \includegraphics[width=\textwidth]{S_1_1_2.png}
        \caption{}
    \end{subfigure}\\[0.5ex]
    \begin{subfigure}[b]{0.42\textwidth}
        \includegraphics[width=\textwidth]{S_1_2_1.png}
        \caption{}
    \end{subfigure}
    \hfill
    \begin{subfigure}[b]{0.42\textwidth}
        \includegraphics[width=\textwidth]{S_1_2_2.png}
        \caption{}
    \end{subfigure}\\[0.5ex]
    \begin{subfigure}[b]{0.42\textwidth}
        \includegraphics[width=\textwidth]{S_1_3_1.png}
        \caption{}
    \end{subfigure}
    \hfill
    \begin{subfigure}[b]{0.42\textwidth}
        \includegraphics[width=\textwidth]{S_1_3_2.png}
        \caption{}
    \end{subfigure}
    \caption{Streamline patterns for lid-driven cavity flow computed using central difference scheme. 
(a,b) $\text{Re} = 100$; (c,d) $\text{Re} = 1000$; (e,f) $\text{Re} = 5000$. 
(a,c,e) $81\times81$ grid; (b,d,f) $161\times161$ grid.}
    \label{fig:Streamline_CDs}
\end{figure}

\subsubsection{QUICK Scheme}
\begin{figure}[H]
    \centering
    \begin{subfigure}[b]{0.42\textwidth}
        \includegraphics[width=\textwidth]{S_2_1_1.png}
        \caption{}
    \end{subfigure}
    \hfill
    \begin{subfigure}[b]{0.42\textwidth}
        \includegraphics[width=\textwidth]{S_2_1_2.png}
        \caption{}
    \end{subfigure}\\[0.5ex]
    \begin{subfigure}[b]{0.42\textwidth}
        \includegraphics[width=\textwidth]{S_2_2_1.png}
        \caption{}
    \end{subfigure}
    \hfill
    \begin{subfigure}[b]{0.42\textwidth}
        \includegraphics[width=\textwidth]{S_2_2_2.png}
        \caption{}
    \end{subfigure}\\[0.5ex]
    \begin{subfigure}[b]{0.42\textwidth}
        \includegraphics[width=\textwidth]{S_2_3_1.png}
        \caption{}
    \end{subfigure}
    \hfill
    \begin{subfigure}[b]{0.42\textwidth}
        \includegraphics[width=\textwidth]{S_2_3_2.png}
        \caption{}
    \end{subfigure}
    \caption{Streamline patterns for lid-driven cavity flow computed using QUICK scheme. 
(a,b) $\text{Re} = 100$; (c,d) $\text{Re} = 1000$; (e,f) $\text{Re} = 5000$. 
(a,c,e) $81\times81$ grid; (b,d,f) $161\times161$ grid.}    
    \label{fig:Streamline_QUICK}
\end{figure}

\subsubsection{TVD-MUSCL Scheme}
\begin{figure}[H]
    \centering
    \begin{subfigure}[b]{0.42\textwidth}
        \includegraphics[width=\textwidth]{S_3_1_1.png}
        \caption{}
    \end{subfigure}
    \hfill
    \begin{subfigure}[b]{0.42\textwidth}
        \includegraphics[width=\textwidth]{S_3_1_2.png}
        \caption{}
    \end{subfigure}\\[0.5ex]
    \begin{subfigure}[b]{0.42\textwidth}
        \includegraphics[width=\textwidth]{S_3_2_1.png}
        \caption{}
    \end{subfigure}
    \hfill
    \begin{subfigure}[b]{0.42\textwidth}
        \includegraphics[width=\textwidth]{S_3_2_2.png}
        \caption{}
    \end{subfigure}\\[0.5ex]
    \begin{subfigure}[b]{0.42\textwidth}
        \includegraphics[width=\textwidth]{S_3_3_1.png}
        \caption{}
    \end{subfigure}
    \hfill
    \begin{subfigure}[b]{0.42\textwidth}
        \includegraphics[width=\textwidth]{S_3_3_2.png}
        \caption{}
    \end{subfigure}
    \caption{Streamline patterns for lid-driven cavity flow computed using TVD-MUSCL scheme.
    (a,b) $\text{Re} = 100$; (c,d) $\text{Re} = 1000$; (e,f) $\text{Re} = 5000$.
    (a,c,e) $81 \times 81$ grid; (b,d,f) $161 \times 161$ grid.}    
    \label{fig:Streamline_TVD-MUSCL}
\end{figure}

\subsection{Velocity Magnitude Contours}
\noindent This section presents velocity magnitude contour plots for the same three numerical schemes across different Reynolds numbers and grid resolutions.
\subsubsection{Arrangement of Velocity Magnitude Contours}
\noindent The below figure shows the arrangement method of my velocity magnitude contour plots results,
split them into 81×81 and 161×161 parts at first, and separate figures into three parts based on 3 types of Re numbers. 
\begin{figure}[H]
  \centering
  \includegraphics[width=0.8\textwidth]{Arrangement2.jpeg}
  \caption{Arrangement method of velocity magnitude contours: organized by grid resolution (81×81 and 161×161) and Reynolds number (Re = 100, 1000, 5000).}
  \label{fig:arrangement2}
\end{figure}
\subsubsection{Central Difference Scheme (CDs)}
\begin{figure}[H]
    \centering
    \begin{subfigure}[b]{0.42\textwidth}
        \includegraphics[width=\textwidth]{C_1_1_1.png}
        \caption{}
    \end{subfigure}
    \hfill
    \begin{subfigure}[b]{0.42\textwidth}
        \includegraphics[width=\textwidth]{C_1_1_2.png}
        \caption{}
    \end{subfigure}\\[0.5ex]
    \begin{subfigure}[b]{0.42\textwidth}
        \includegraphics[width=\textwidth]{C_1_2_1.png}
        \caption{}
    \end{subfigure}
    \hfill
    \begin{subfigure}[b]{0.42\textwidth}
        \includegraphics[width=\textwidth]{C_1_2_2.png}
        \caption{}
    \end{subfigure}\\[0.5ex]
    \begin{subfigure}[b]{0.42\textwidth}
        \includegraphics[width=\textwidth]{C_1_3_1.png}
        \caption{}
    \end{subfigure}
    \hfill
    \begin{subfigure}[b]{0.42\textwidth}
        \includegraphics[width=\textwidth]{C_1_3_2.png}
        \caption{}
    \end{subfigure}
    \caption{Velocity magnitude contours for lid-driven cavity flow computed using central difference scheme. 
(a,b) $\text{Re} = 100$; (c,d) $\text{Re} = 1000$; (e,f) $\text{Re} = 5000$. 
(a,c,e) $81\times81$ grid; (b,d,f) $161\times161$ grid.}
    \label{fig:Contour_CDs}
\end{figure}

\subsubsection{QUICK Scheme}
\begin{figure}[H]
    \centering
    \begin{subfigure}[b]{0.42\textwidth}
        \includegraphics[width=\textwidth]{C_2_1_1.png}
        \caption{}
    \end{subfigure}
    \hfill
    \begin{subfigure}[b]{0.42\textwidth}
        \includegraphics[width=\textwidth]{C_2_1_2.png}
        \caption{}
    \end{subfigure}\\[0.5ex]
    \begin{subfigure}[b]{0.42\textwidth}
        \includegraphics[width=\textwidth]{C_2_2_1.png}
        \caption{}
    \end{subfigure}
    \hfill
    \begin{subfigure}[b]{0.42\textwidth}
        \includegraphics[width=\textwidth]{C_2_2_2.png}
        \caption{}
    \end{subfigure}\\[0.5ex]
    \begin{subfigure}[b]{0.42\textwidth}
        \includegraphics[width=\textwidth]{C_2_3_1.png}
        \caption{}
    \end{subfigure}
    \hfill
    \begin{subfigure}[b]{0.42\textwidth}
        \includegraphics[width=\textwidth]{C_2_3_2.png}
        \caption{}
    \end{subfigure}
    \caption{Velocity magnitude contours for lid-driven cavity flow computed using QUICK scheme. 
(a,b) $\text{Re} = 100$; (c,d) $\text{Re} = 1000$; (e,f) $\text{Re} = 5000$. 
(a,c,e) $81\times81$ grid; (b,d,f) $161\times161$ grid.}
    \label{fig:Contour_QUICK}
\end{figure}

\subsubsection{TVD-MUSCL Scheme}
\begin{figure}[H]
    \centering
    \begin{subfigure}[b]{0.42\textwidth}
        \includegraphics[width=\textwidth]{C_3_1_1.png}
        \caption{}
    \end{subfigure}
    \hfill
    \begin{subfigure}[b]{0.42\textwidth}
        \includegraphics[width=\textwidth]{C_3_1_2.png}
        \caption{}
    \end{subfigure}\\[0.5ex]
    \begin{subfigure}[b]{0.42\textwidth}
        \includegraphics[width=\textwidth]{C_3_2_1.png}
        \caption{}
    \end{subfigure}
    \hfill
    \begin{subfigure}[b]{0.42\textwidth}
        \includegraphics[width=\textwidth]{C_3_2_2.png}
        \caption{}
    \end{subfigure}\\[0.5ex]
    \begin{subfigure}[b]{0.42\textwidth}
        \includegraphics[width=\textwidth]{C_3_3_1.png}
        \caption{}
    \end{subfigure}
    \hfill
    \begin{subfigure}[b]{0.42\textwidth}
        \includegraphics[width=\textwidth]{C_3_3_2.png}
        \caption{}
    \end{subfigure}
    \caption{Velocity magnitude contours for lid-driven cavity flow computed using TVD-MUSCL scheme. 
(a,b) $\text{Re} = 100$; (c,d) $\text{Re} = 1000$; (e,f) $\text{Re} = 5000$. 
(a,c,e) $81\times81$ grid; (b,d,f) $161\times161$ grid.}
    \label{fig:Contour_MUSCL}
\end{figure}

\subsection{Velocity Centerline Comparison for Different Re Numbers}
\noindent This section presents the centerline velocity profiles comparing different Reynolds numbers for each numerical scheme.
\subsubsection{Arrangement of u-Velocity and v-Velocity Profiles}
\noindent Each subfigure presents results for three Reynolds numbers ($\text{Re} = 100, 1000, 5000$). The arrangement method is summarized below: 
(1) $u$-velocity profiles on the right-hand side, and $v$-velocity profiles on the left-hand side; 
(2) Central Difference Scheme at top, QUICK in the middle, and TVD-MUSCL at bottom.
\begin{figure}[H]
  \centering
  \includegraphics[width=0.7\textwidth]{Arrangement3.jpeg}
  \caption{Arrangement method of velocity centerline profiles. Each row corresponds to a numerical scheme. Each column shows the $u$-velocity profiles on the left-hand side and the $v$-velocity profiles on the right-hand side. Results are presented for three Reynolds numbers ($\text{Re} = 100, 1000, 5000$) and two grid resolutions ($81 \times 81$ and $161 \times 161$).}
  \label{fig:arrangement3}
\end{figure}

\subsubsection{Centerline Velocity Profiles - Grid 81x81}
\begin{figure}[H]
    \centering
    \begin{subfigure}[b]{0.48\textwidth}
        \includegraphics[width=\textwidth]{V_1_1_1.png}
        \caption{}
    \end{subfigure}
    \hfill
    \begin{subfigure}[b]{0.46\textwidth}
        \includegraphics[width=\textwidth]{V_1_1_2.png}
        \caption{}
    \end{subfigure}\\[0.5ex]
    \begin{subfigure}[b]{0.48\textwidth}
        \includegraphics[width=\textwidth]{V_1_2_1.png}
        \caption{}
    \end{subfigure}
    \hfill
    \begin{subfigure}[b]{0.46\textwidth}
        \includegraphics[width=\textwidth]{V_1_2_2.png}
        \caption{}
    \end{subfigure}\\[0.5ex]
    \begin{subfigure}[b]{0.48\textwidth}
        \includegraphics[width=\textwidth]{V_1_3_1.png}
        \caption{}
    \end{subfigure}
    \hfill
    \begin{subfigure}[b]{0.46\textwidth}
        \includegraphics[width=\textwidth]{V_1_3_2.png}
        \caption{}
    \end{subfigure}
    \caption{Velocity profiles for lid-driven cavity flow at $\text{Re} = 1000$ on $81\times81$ grid. 
(a,c,e) $u$-velocity along vertical centreline; (b,d,f) $v$-velocity along horizontal centreline. 
(a,b) Central difference scheme; (c,d) QUICK scheme; (e,f) TVD-MUSCL scheme. 
Symbols: benchmark data from Ghia \textit{et al.}~(1982).}
    \label{fig:Velocity_Centerline_Grid 81x81}
\end{figure}

\subsubsection{Centerline Velocity Profiles - Grid 161x161}
\begin{figure}[H]
    \centering
    \begin{subfigure}[b]{0.48\textwidth}
        \includegraphics[width=\textwidth]{V_2_1_1.png}
        \caption{}
    \end{subfigure}
    \hfill
    \begin{subfigure}[b]{0.46\textwidth}
        \includegraphics[width=\textwidth]{V_2_1_2.png}
        \caption{}
    \end{subfigure}\\[0.5ex]
    \begin{subfigure}[b]{0.48\textwidth}
        \includegraphics[width=\textwidth]{V_2_2_1.png}
        \caption{}
    \end{subfigure}
    \hfill
    \begin{subfigure}[b]{0.46\textwidth}
        \includegraphics[width=\textwidth]{V_2_2_2.png}
        \caption{}
    \end{subfigure}\\[0.5ex]
    \begin{subfigure}[b]{0.48\textwidth}
        \includegraphics[width=\textwidth]{V_2_3_1.png}
        \caption{}
    \end{subfigure}
    \hfill
    \begin{subfigure}[b]{0.46\textwidth}
        \includegraphics[width=\textwidth]{V_2_3_2.png}
        \caption{}
    \end{subfigure}
    \caption{Velocity profiles for lid-driven cavity flow at $\text{Re} = 1000$ on $161\times161$ grid. 
(a,c,e) $u$-velocity along vertical centreline; (b,d,f) $v$-velocity along horizontal centreline. 
(a,b) Central difference scheme; (c,d) QUICK scheme; (e,f) TVD-MUSCL scheme. 
Symbols: benchmark data from Ghia \textit{et al.}~(1982).}
    \label{fig:Velocity_Centerline_Grid 161x161}
\end{figure}


\subsection{Velocity Centerline Comparison for Different Schemes}
\begin{figure}[H]
    \centering
    \begin{subfigure}[b]{0.48\textwidth}
        \includegraphics[width=\textwidth]{V_1_4.png}
        \caption{}
    \end{subfigure}
    \hfill
    \begin{subfigure}[b]{0.48\textwidth}
        \includegraphics[width=\textwidth]{V_2_4.png}
        \caption{}
    \end{subfigure}\\[0.5ex]
    
    \begin{subfigure}[b]{0.48\textwidth}
        \includegraphics[width=\textwidth]{V_1_5.png}
        \caption{}
    \end{subfigure}
    \hfill
    \begin{subfigure}[b]{0.48\textwidth}
        \includegraphics[width=\textwidth]{V_2_5.png}
        \caption{}
    \end{subfigure}\\[0.5ex]
    
    \begin{subfigure}[b]{0.48\textwidth}
        \includegraphics[width=\textwidth]{V_1_6.png}
        \caption{}
    \end{subfigure}
    \hfill
    \begin{subfigure}[b]{0.48\textwidth}
        \includegraphics[width=\textwidth]{V_2_6.png}
        \caption{}
    \end{subfigure}
    
   \caption{Centreline velocity profiles for different numerical schemes at $\text{Re} = 100$ (a,b), $1000$ (c,d), and $5000$ (e,f). 
Left column: $81\times81$ grid; right column: $161\times161$ grid. 
Each panel shows $u$-velocity (left) and $v$-velocity (right) along centrelines. 
Symbols: Ghia \textit{et al.}~(1982).}
    \label{fig:velocity_comparison_schemes}
\end{figure}

\subsection{Velocity Centerline Comparison for Different Grid Sizes at High Re Number}
\begin{figure}[htbp]
    \centering
    \begin{subfigure}[b]{0.48\textwidth}
        \includegraphics[width=\textwidth]{V_3_1.png}
        \caption{}
        \label{fig:grid_conv_cds}
    \end{subfigure}
    \hfill
    \begin{subfigure}[b]{0.48\textwidth}
        \includegraphics[width=\textwidth]{V_3_2.png}
        \caption{}
        \label{fig:grid_conv_quick}
    \end{subfigure}\\[1ex]
    
    \begin{subfigure}[b]{0.48\textwidth}
        \includegraphics[width=\textwidth]{V_3_3.png}
        \caption{}
        \label{fig:grid_conv_tvd}
    \end{subfigure}
    
    \caption{Grid convergence study for $u$-velocity profiles along vertical centreline at $\text{Re} = 5000$. 
    (a) Central difference scheme; (b) QUICK scheme; (c) TVD-MUSCL scheme. 
    Four grid resolutions: $81\times81$, $113\times113$, $161\times161$, and $227\times227$. 
    Symbols: benchmark data from Ghia \textit{et al.}~(1982).}
    \label{fig:grid_convergence_re5000}
\end{figure}
\noindent Whatever discrete scheme I used, from figure \ref{fig:grid_convergence_re5000} we can see that when grid size is smaller(like $227\times227$), centerline velocity profiles for $u$-Velocity and $v$-Velocity are closer to the benchmark by Ghia\textrm{et al.}(1982).


\newpage 
\subsection{Accuracy Comparison With Different Schemes}

For grid size $81 \times 81$, $113 \times 113$, $161 \times 161$, and $227 \times 227$,
I have completed the evaluation of the normalized $L_{1}$ errors for both the centerline 
$u$- and 
$v$-velocity profiles at a Reynolds number of 5000 across all grid sizes, using Ghia’s benchmark data as the exact solution.
\begin{table}[H]
\centering
\small
\caption{$L_1$ error and grid convergence data for different numerical schemes. 
The $L_1$ norm is computed as $L_1 \text{ norm} = \frac{1}{34}\sum_{i=1}^{17}\sum_{j=1}^{17}(|u-u_{\text{Ghia}}|+|v-v_{\text{Ghia}}|)$.}
\label{tab:error_comparison}
\begin{tabular*}{\textwidth}{@{\extracolsep{\fill}}lcccccccc@{}}
\toprule
\multirow{2}{*}{Grid size} & \multirow{2}{*}{$\ln(\Delta x)$} & \multicolumn{2}{c}{CDs} & \multicolumn{2}{c}{QUICK} & \multicolumn{2}{c}{TVD-MUSCL} \\
\cmidrule(lr){3-4} \cmidrule(lr){5-6} \cmidrule(lr){7-8}
& & $L_1$ Error & $\ln(L_1)$ & $L_1$ Error & $\ln(L_1)$ & $L_1$ Error & $\ln(L_1)$ \\
\midrule
$81 \times 81$   & $-4.382$ & 1.79073  & $-2.943$ & 1.26386  & $-3.289$ & 1.05693  & $-3.464$ \\
$113 \times 113$ & $-4.718$ & 0.924517 & $-3.686$ & 0.593034 & $-4.149$ & 0.744874 & $-4.129$ \\
$161 \times 161$ & $-5.075$ & 0.510966 & $-4.280$ & 0.467266 & $-4.374$ & 0.534575 & $-4.246$ \\
$227 \times 227$ & $-5.421$ & 0.279889 & $-4.882$ & 0.305895 & $-4.794$ & 0.370451 & $-4.615$ \\
\bottomrule
\end{tabular*}
\end{table}

\begin{figure}[H]
\centering
\begin{tikzpicture}
\begin{axis}[
    width=0.7\textwidth,
    height=0.55\textwidth,
    xlabel={$\ln(\Delta x)$},
    ylabel={$\ln(L_1 \text{ norm})$},
    title={Accuracy Comparison},
    title style={font=\bfseries, at={(0.5,1.05)}},
    legend pos=south east,
    legend style={
        font=\small, 
        draw=black, 
        fill=white,
        cells={anchor=west}
    },
    grid=major,
    grid style={line width=0.1pt, draw=gray!30},
    xmin=-5.5, xmax=-4.2,
    ymin=-5.1, ymax=-2.8,
    tick label style={font=\small},
    label style={font=\normalsize},
    mark size=4pt,
]

% CDs - 中空圓圈(紅色) + 回歸線
\addplot[
    only marks, 
    mark=o, 
    mark options={draw=red, fill=white, line width=1pt}
] coordinates {
    (-4.382, -2.943)
    (-4.718, -3.686)
    (-5.075, -4.280)
    (-5.421, -4.882)
};

\addplot[red, thick, domain=-5.5:-4.3, samples=2, forget plot] 
{1.77375*x + 4.833};

\addlegendentry{CDs (slope=1.77375)}

% QUICK - 中空圓圈(藍色) + 回歸線
\addplot[
    only marks, 
    mark=o, 
    mark options={draw=blue, fill=white, line width=1pt}
] coordinates {
    (-4.382, -3.289)
    (-4.718, -4.149)
    (-5.075, -4.374)
    (-5.421, -4.794)
};

\addplot[blue, thick, domain=-5.5:-4.3, samples=2, forget plot] 
{1.29119*x + 2.368};

\addlegendentry{QUICK (slope=1.29119)}

% TVD-MUSCL - 中空三角形(綠色) + 回歸線
\addplot[
    only marks, 
    mark=triangle, 
    mark options={draw=green!60!black, fill=white, line width=1pt},
    mark size=5pt
] coordinates {
    (-4.382, -3.464)
    (-4.718, -4.129)
    (-5.075, -4.246)
    (-5.421, -4.615)
};

\addplot[green!60!black, thick, domain=-5.5:-4.3, samples=2, forget plot] 
{1.08102*x + 1.270};

\addlegendentry{TVD-MUSCL (slope=1.08102)}

% 添加公式註解(可選)
\node[anchor=north west, font=\footnotesize] at (axis cs:-5.4,-3.0) 
{$L_1 norm = \frac{1}{34}\left(\sum_{i=1}^{17}\sum_{j=1}^{17}(|u-u_{Ghia}|+|v-v_{Ghia}|)\right)$};

\end{axis}
\end{tikzpicture}
\caption{Accuracy comparison for difference schemes based on $L_{1}$ norm at Re $=5000$. As the data shown in Table \ref{tab:error_comparison}, the slope of the line decided by $ln(L_{1}nrom)$ and $ln(dx)$ is the accuracy of the scheme. The accuracy for CDs, QUICK, and TVD-MUSCL are approximately 1.77375, 1.29119. and 1.08102, respectively.}
\label{fig:accuracy}
\end{figure}
In the 227×227 grid, the error compared with the benchmark for CDS is less than that for  QUICK.This may be due to the limited number of benchmark data points.The numerical solutions in the figure represent all my experience with various relaxation factor sets. The unusual finding is that inappropriate relaxation settings lead to either incorrect solutions or results that do not match the benchmark, so I maybe have a chance to find diference solution with different set.

%% ============================================
%% Appendix: Source Code
%% ============================================
\newpage
\appendix
% Appendix 設定：章節以 A, B, C 排列，子章節以 A.1, A.2 排列
\renewcommand{\thesection}{\Alph{section}}
\renewcommand{\thesubsection}{\thesection.\arabic{subsection}}
\renewcommand{\thesubsubsection}{\thesubsection.\arabic{subsubsection}}
\renewcommand{\theequation}{\thesection.\arabic{equation}}
\renewcommand{\thefigure}{\thesection.\arabic{figure}}
\renewcommand{\thetable}{\thesection.\arabic{table}}

\section{Source Code}
\label{app:source_code}

\subsection{Header Files (.h files)}

\subsubsection{TDMA Solver}
\label{subsec:code1} 
\lstinputlisting[language=C++, caption={0.TDMA\_solver.h}]{0.TDMA_solver.h}
%\textit{(Code omitted for faster compilation - uncomment when finalizing)}

\subsubsection{Pressure Correction Equation}
\label{subsec:code2} 
\lstinputlisting[language=C++, caption={0.P\_correction\_eq.h}]{0.P_correction_eq.h}
%\textit{(Code omitted for faster compilation - uncomment when finalizing)}

\subsubsection{Pressure Correction Equation with TDMA}
\label{subsec:code3} 
\lstinputlisting[language=C++, caption={0.P\_correction\_eq\_TDMA.h}]{0.P_correction_eq_TDMA.h}
%\textit{(Code omitted for faster compilation - uncomment when finalizing)}

\subsection{Main Programs (.cpp files)}

\subsubsection{FDM Central Difference Scheme}
\label{subsec:code4} 
\lstinputlisting[language=C++, caption={1.FDM\_CDs.cpp}]{1.FDM_CDs.cpp}
%\textit{(Code omitted for faster compilation - uncomment when finalizing)}

\subsubsection{FVM Non-Deferred Correction}
\label{subsec:code5} 
\lstinputlisting[language=C++, caption={2.FVM\_non-deferred correction.cpp}]{2.FVM_non-deferred correction.cpp}
%\textit{(Code omitted for faster compilation - uncomment when finalizing)}

\subsubsection{FVM Deferred Correction}
\label{subsec:code6} 
\lstinputlisting[language=C++, caption={3.FVM\_deferred correction.cpp}]{3.FVM_deferred correction.cpp}
%\textit{(Code omitted for faster compilation - uncomment when finalizing)}

\subsubsection{FVM Deferred Correction with ADI+TDMA}
\label{subsec:code7} 
\lstinputlisting[language=C++, caption={4.FVM\_deferred correction\_ADI+TDMA.cpp}]{4.FVM_deferred correction_ADI+TDMA.cpp}
%\textit{(Code omitted for faster compilation - uncomment when finalizing)}

%% ============================================
%% Appendix B: SIMPLE Algebraic System
%% ============================================
\newpage
\section{SIMPLE Algebraic System}
\label{app:simple_algebraic}

\begin{table}[H]
\centering
\begin{tabular}{|c|c|}
\hline
\rowcolor{orange!80}
\textbf{Step} & \textbf{Equation} \\
\hline
1 
& 
\begin{minipage}[c]{0.85\textwidth}
\centering
\vspace{0.3em}
\textbf{predicted form of the discrete momentum integration equation\eqref{eq:eq1}}
\vspace{0.3em}
\end{minipage} \\
\hline
2 & 
\begin{minipage}[c]{0.85\textwidth}
\centering
\vspace{0.3em}
\textbf{pressure correction equation (pressure Poisson equation)\eqref{eq:eq2}}
\vspace{0.3em}
\end{minipage} \\
\hline
3 & 
\begin{minipage}[c]{0.85\textwidth}
\centering
\vspace{0.3em}
\textbf{pressure predicted-correction relationship\eqref{eq:eq3}}
\vspace{0.3em}
\end{minipage} \\
\hline
4 & 
\begin{minipage}[c]{0.85\textwidth}
\centering
\vspace{0.3em}
\textbf{Patankar's correction form of the discrete momentum integration equation\eqref{eq:eq4}}
\vspace{0.3em}
\end{minipage} \\
\hline
5 & 
\begin{minipage}[c]{0.85\textwidth}
\centering
\vspace{0.3em}
\textbf{velocity predicted-correction relationship\eqref{eq:eq5}}
\vspace{0.3em}
\end{minipage} \\
\hline
6 & 
\begin{minipage}[c]{0.85\textwidth}
\centering
\vspace{0.3em}
\textbf{predicted form of the continuity integration equation\eqref{eq:eq6}}
\vspace{0.3em}
\end{minipage} \\
\hline
7 & 
\begin{minipage}[c]{0.85\textwidth}
\centering
\vspace{0.3em}
\textbf{updating predicted pressure from correct pressure\eqref{eq:eq7}}
\vspace{0.3em}
\end{minipage} \\
\hline
8 & 
\begin{minipage}[c]{0.85\textwidth}
\centering
\vspace{0.3em}
\textbf{updating predicted velocity from correct velocity\eqref{eq:eq8}}
\vspace{0.3em}
\end{minipage} \\
\hline
\end{tabular}
\caption{SIMPLE algorithm discrete equation system}
\label{tab:SIMPLE_equations}
\end{table}


This appendix provides the detailed discrete equations corresponding to the
steps summarised in Table~B.1. The formulations follow the standard
SIMPLE procedure for pressure–velocity coupling on a staggered grid.

\renewcommand\theequation{B.\arabic{equation}}
\setcounter{equation}{0}

% ---------------------------------------------------
\subsection*{B.1 \quad Step 1: Predicted Momentum Equation}

The discrete $x$–momentum equation for the predicted velocity
$u_{i,J}^{*,\,N+1/2}$ is written as
\begin{equation}
\begin{aligned}
(a_P)_{i,J}^{*,\,N}\,u_{i,J}^{*,\,N+1/2}
&\;+\;(a_{WW})_{i,J}^{*,\,N} u_{i-2,J}^{*,\,N}
+\;(a_W)_{i,J}^{*,\,N}   u_{i-1,J}^{*,\,N}
+\;(a_E)_{i,J}^{*,\,N}   u_{i+1,J}^{*,\,N}
+\;(a_{EE})_{i,J}^{*,\,N} u_{i+2,J}^{*,\,N} \\
&\;+\;(a_{SS})_{i,J}^{*,\,N} u_{i,J-2}^{*,\,N}
+\;(a_S)_{i,J}^{*,\,N}   u_{i,J-1}^{*,\,N}
+\;(a_N)_{i,J}^{*,\,N}   u_{i,J+1}^{*,\,N}
+\;(a_{NN})_{i,J}^{*,\,N} u_{i,J+2}^{*,\,N}
\\[3pt]
&= -\,\frac{ p_{I,J}^{*,\,N} - p_{I-1,J}^{*,\,N} }{\rho}\,\delta y.
\end{aligned}
\label{eq:eq1}
\end{equation}
% ---------------------------------------------------

\subsection*{B.2 \quad Step 2: Pressure-Correction Equation}

Using the discrete continuity constraint, the pressure–correction equation becomes
\begin{equation}
\begin{aligned}
&\Bigg(
\frac{\delta y}{\rho(a_P)_{i+1,J}^{*,\,N+1/2}}
+ \frac{\delta y}{\rho(a_P)_{i,J}^{*,\,N+1/2}}
+ \frac{\delta x}{\rho(b_P)_{I,j+1}^{*,\,N+1/2}}
+ \frac{\delta x}{\rho(b_P)_{I,j}^{*,\,N+1/2}}
\Bigg)
p_{I,J}^{\prime\,N+1}
\\[4pt]
&= 
\frac{\delta x}{\rho(b_P)_{I,j}^{*,\,N+1/2}}\, p_{I,J-1}^{\prime\,N}
\;+\;
\frac{\delta x}{\rho(b_P)_{I,j+1}^{*,\,N+1/2}}\, p_{I,J+1}^{\prime\,N}
\\
&\quad
+\frac{\delta y}{\rho(a_P)_{i,J}^{*,\,N+1/2}} \, p_{I-1,J}^{\prime\,N}
\;+\;
\frac{\delta y}{\rho(a_P)_{i+1,J}^{*,\,N+1/2}} \, p_{I+1,J}^{\prime\,N}
\\
&\quad
+\;u_{i,J}^{*,\,N+1/2} - u_{i+1,J}^{*,\,N+1/2}
\;+\;
v_{I,j}^{*,\,N+1/2} - v_{I,j+1}^{*,\,N+1/2}.
\end{aligned}
\label{eq:eq2}
\end{equation}
% ---------------------------------------------------

\subsection*{B.3 \quad Step 3: Pressure Update}

The corrected pressure is updated using the relaxation factor $\alpha_p$:
\begin{equation}
p_{I,J}^{N+1}
= p_{I,J}^{*,\,N}
+ \alpha_p\, p_{I,J}^{\prime\,N+1}.
\label{eq:eq3}
\end{equation}

% ---------------------------------------------------
\subsection*{B.4 \quad Step 4: Velocity Correction}

The velocity correction follows directly from the linearised momentum operator:
\begin{equation}
u_{i,J}^{\prime\,N+1}
= -\,\frac{ p_{I,J}^{\prime\,N+1} - p_{I-1,J}^{\prime\,N+1} }
       {\rho\, (a_P)_{i,J}^{*,\,N+1/2}}\,
  \delta y.
\label{eq:eq4}
\end{equation}

% ---------------------------------------------------
\subsection*{B.5 \quad Step 5: Velocity Update}

\begin{equation}
u_{i,J}^{N+1}
= u_{i,J}^{*,\,N+1/2}
+ u_{i,J}^{\prime\,N+1}.
\label{eq:eq5}
\end{equation}

% ---------------------------------------------------
\subsection*{B.6 \quad Step 6: Continuity Check}

A discrete conservation check is performed:
\begin{equation}
A_e\,u_{i+1,J}^{*,\,N+1/2}
- A_w\,u_{i,J}^{*,\,N+1/2}
+ A_n\,v_{I,j+1}^{*,\,N+1/2}
- A_s\,v_{I,j}^{*,\,N+1/2}
= 0.
\label{eq:eq6}
\end{equation}

% ---------------------------------------------------
\subsection*{B.7–B.8 \quad Step 7–8: Update of Predicted Fields}


\begin{equation}
\label{eq:eq7}
p_{I,J}^{*,\,N+1} = p_{I,J}^{N+1}.
\end{equation}
\begin{equation}
\label{eq:eq8}
u_{i,J}^{*,\,N+1} = u_{i,J}^{N+1}
\end{equation}

\end{document}